
\section*{Round 2 continuation for Erd\H{o}s Problems \#278, \#279, \#281}

We write $d(S)$ for natural density when it exists, $\overline d(S)$ for upper density, and $\underline d(S)$ for lower density.

\section*{Problem 281}

\subsection*{1) ROUND-2 OBJECTIVE}

We pursue \textbf{(A) proof}. Round 1 already gave a full proof using profinite/Haar/Dini machinery. The most promising Round-2 task is therefore not to restart from scratch, but to close the argument more cleanly by replacing that longer proof with a short classical proof built directly on the finite-prefix extremal statement already used in Round 1.

\subsection*{2) Round-1 FOUNDATION USED}

We use two specific Round-1 facts.

\begin{enumerate}
\item For a fixed finite set of moduli $n_1,\dots,n_k$, the density of integers \emph{avoiding} chosen congruence classes is maximised when the residues are aligned (equivalently, all equal modulo their moduli). This is the finite-prefix extremal theorem from Simpson already used in Round 1.
\item For aligned residues $a_i\equiv 0$, the avoided set for the first $k$ moduli is exactly the set of integers divisible by none of $n_1,\dots,n_k$.
\end{enumerate}

\subsection*{3) NEW INSIGHT / TOOL (ROUND-2)}

The new tool is a short combination of:

\begin{enumerate}
\item the Davenport--Erd\H{o}s theorem on sets of multiples, and
\item the Round-1 finite-prefix extremal inequality.
\end{enumerate}

This removes the need for the more elaborate profinite/Dini argument and yields a very short proof.

\subsection*{4) ATTACK PLAN (ROUND-2)}

The only real gap left after Round 1 was conceptual economy: can the conclusion be proved directly from classical density theorems?

Yes. The plan is:

\begin{enumerate}
\item specialise the hypothesis to the zero residues;
\item apply Davenport--Erd\H{o}s to the set of multiples of $\{n_i\}$;
\item use the finite-prefix extremal inequality from Round 1 to transfer the estimate from the zero residues to arbitrary residues.
\end{enumerate}

\subsection*{5) WORK (ROUND-2)}

Let
\[
M:=\{\,m\in\mathbf N : n_i\mid m \text{ for at least one } i\ge 1\,\}
\]
be the set of multiples of the sequence $(n_i)$, and let
\[
M_k:=\{\,m\in\mathbf N : n_i\mid m \text{ for at least one } 1\le i\le k\,\}.
\]
Also let
\[
F_k(\mathbf a):=\{\,m\in\mathbf Z : m\not\equiv a_i\pmod{n_i}\ \text{for all } 1\le i\le k\,\}.
\]

By hypothesis, for every residue choice $(a_i)$ the full avoidance set has density $0$. In particular, for the special choice $a_i\equiv 0\pmod{n_i}$ for all $i$, the set of integers avoiding all these zero congruences has density $0$. But that set is exactly the complement of $M$. Hence
\[
d(M)=1.
\]

Now invoke the Davenport--Erd\H{o}s theorem for sets of multiples: for an increasing sequence $(n_i)$,
\[
\lim_{k\to\infty} d(M_k)=\underline d(M).
\]
Since $d(M)=1$, we have $\underline d(M)=1$, and therefore
\[
\lim_{k\to\infty} d(M_k)=1.
\]
Equivalently,
\[
\lim_{k\to\infty} d(F_k(\mathbf 0))=0,
\]
because $F_k(\mathbf 0)$ is the complement of $M_k$.

Now fix $k$. By the finite-prefix extremal theorem from Round 1 (Simpson's theorem), for every choice of residues $(a_i)$,
\[
d(F_k(\mathbf a))\le d(F_k(\mathbf 0)).
\]
Since the right-hand side tends to $0$ as $k\to\infty$, given any $\epsilon>0$ we may choose $k$ so large that
\[
d(F_k(\mathbf 0))<\epsilon.
\]
Then automatically, for every residue choice $(a_i)$,
\[
d(F_k(\mathbf a))\le d(F_k(\mathbf 0))<\epsilon.
\]

This is exactly the required conclusion.

\hfill $\square$

\subsection*{6) ADVERSARIAL VERIFICATION}

\begin{enumerate}
\item \textbf{Did we use the hypothesis for all residue choices?} Yes, but only once: we specialised it to the zero choice $a_i\equiv 0$ to deduce $d(M)=1$.

\item \textbf{Are the hypotheses of Davenport--Erd\H{o}s satisfied?} Yes. The theorem applies to an increasing sequence of positive integers and the associated set of multiples.

\item \textbf{Where is the finite-prefix extremal theorem used?} Exactly in the inequality
\[
d(F_k(\mathbf a))\le d(F_k(\mathbf 0)).
\]
Without that step, information about zero residues would not automatically transfer to arbitrary residues.

\item \textbf{Edge case $n_1=1$.} Then $M=\mathbf N$ from the start, $d(M_k)=1$ for all $k\ge 1$, and the conclusion is immediate.
\end{enumerate}

\subsection*{7) FINAL}

\textbf{FULL SOLUTION --- FULL PROOF.}

Round 2 upgrades the Round-1 solution by giving a short classical proof: the statement follows directly from the Davenport--Erd\H{o}s theorem on sets of multiples together with the finite-prefix extremal inequality from Simpson.

\subsection*{8) COMPLETION ESTIMATE}

COMPLETION: 100\%.

\subsection*{9) REFERENCES}

\begin{enumerate}
\item P. Erd\H{o}s and R. L. Graham, \emph{Old and New Problems and Results in Combinatorial Number Theory}, Monographies de L'Enseignement Math\'ematique 28, Geneva, 1980.
\item H. Davenport and P. Erd\H{o}s, \emph{On sequences of positive integers}, Acta Arithmetica 2 (1936), 147--151.
\item R. J. Simpson, \emph{Exact coverings of the integers by arithmetic progressions}, Discrete Mathematics 59 (1986), 181--190.
\end{enumerate}
